%% ma: L11-B01-0005
%% lop: 11
%% chuong: 1
%% bai: 1
%% dang: 11.1.5
%% loai: TLN
%% muc: TH
%% dap_an: "-0,4"
%% nguon: "de-cuong-giua-ky-1-toan-11-nam-2026-2027-truong-thpt-hoang-van-thu-ha-noi-2-8.pdf (M08, Đề cương giữa kì 1 Toán 11 THPT Hoàng Văn Thụ, Hà Nội), Phần 3 Câu 3"
%% kien_thuc: "Bài 1 (giá trị lượng giác của góc lượng giác)"
%% trang_thai: chinh-thuc
%% ghi_chu: "Dữ kiện α ∈ (−π/2; 0) không cần cho lời giải (\"hàng phần chục\" hiểu là hàng phần mười). Đã duyệt 29/10/2026."
\begin{ex}
Cho $\sin\alpha+\cos\alpha=\dfrac13$, $\alpha\in\left(-\dfrac{\pi}{2};0\right)$. Tính giá trị biểu thức $A=\sin\alpha\cdot\cos\alpha$ (kết quả làm tròn đến hàng phần chục).
\shortans{-0,4}
\loigiai{$(\sin\alpha+\cos\alpha)^2=1+2\sin\alpha\cos\alpha=\dfrac19\Rightarrow A=-\dfrac49\approx-0{,}4$.}
\end{ex}
